\section*{第 \ref{ch:b1:crystals-ionic-covalent} 章——晶体 II：离子、共价与分子固体}

\begin{solution}{exo:b1:crystals-ionic-covalent:1}
阴离子：$8 \times \frac18 + 6 \times \frac12 = 4$；阳离子：$12 \times
\frac14 + 1 = 4$。每个离子有 6 个异号近邻（6:6）。
四个阳离子和四个阴离子：\ce{NaCl}。
\end{solution}

\begin{solution}{exo:b1:crystals-ionic-covalent:2}
体对角线的一半：$a\sqrt3/2 = 412.3 \times 0.866 = \qty{357.1}{pm}$。
一个 \ce{Cs+} 和 $8 \times \frac18 = 1$ 个 \ce{Cl-}，即一个化学式单位。
\end{solution}

\begin{solution}{exo:b1:crystals-ionic-covalent:3}
$x = 102/181 = 0.56$，介于 0.414 与 0.732 之间：八面体配位
（6），属于氯化钠型。
\end{solution}

\begin{solution}{exo:b1:crystals-ionic-covalent:4}
金属：铜。离子型：氯化钠、氟化钙。共价型：
金刚石、石英。分子型：碘、冰。
\end{solution}

\begin{solution}{exo:b1:crystals-ionic-covalent:5}
沿棱方向，阴离子–阳离子–阴离子相接触：$a = 2(r_+ + r_-)$。沿
面对角线有两个阴离子，距离为 $a\sqrt2/2$，二者不得重叠：
$a\sqrt2/2 \ge 2r_-$，于是 $(r_+ + r_-)\sqrt2 \ge 2r_-$，即
$r_+/r_- \ge \sqrt2 - 1$。
\end{solution}

\begin{solution}{exo:b1:crystals-ionic-covalent:6}
$\rho = 4 \times 58.44/(6.022 \times 10^{23} \times (5.641 \times
10^{-8})^3) = \qty{2.16}{g/cm^3}$，对比实测值 \qty{2.17}{g/cm^3}。
\end{solution}

\begin{solution}{exo:b1:crystals-ionic-covalent:7}
$\rho = 168.36/(6.022 \times 10^{23} \times (4.123 \times 10^{-8})^3) =
\qty{3.99}{g/cm^3}$，即实测值。
\end{solution}

\begin{solution}{exo:b1:crystals-ionic-covalent:8}
\ce{Zn-S} $= a\sqrt3/4 = \qty{234.2}{pm}$;
$\rho = 4 \times 97.44/(6.022 \times 10^{23} \times (5.409 \times
10^{-8})^3) = \qty{4.09}{g/cm^3}$，接近
天然闪锌矿的实测值 \qty{4.04}{g/cm^3}。
\end{solution}

\begin{solution}{exo:b1:crystals-ionic-covalent:9}
\ce{C-C} $= 245.6/\sqrt3 = \qty{141.8}{pm}$；层间距 $669.6/2 =
\qty{334.8}{pm}$。晶胞体积 $\frac{\sqrt3}{2}a^2c = \qty{3.498e-23}{cm^3}$，
$\rho = 4 \times 12.01/(6.022 \times 10^{23} \times 3.498 \times 10^{-23})
= \qty{2.28}{g/cm^3}$，对比金刚石的 \qty{3.52}{g/cm^3}：层与层
之间相距甚远，仅由伦敦色散力维系。
\end{solution}

\begin{solution}{exo:b1:crystals-ionic-covalent:10}
见\cref{prop:b1:crystals-ionic-covalent:diamond}的证明：
$C = \pi\sqrt3/16 = 0.34$。硬度来自共价键的强度及其三维刚性，
而非堆积方式：每个原子只有四个近邻，
键角固定，
因而留有大量空隙。
\end{solution}

\begin{solution}{exo:b1:crystals-ionic-covalent:11}
\ce{Si-Si} $= a\sqrt3/4 = \qty{235.2}{pm}$，共价半径为
\qty{117.6}{pm}; $\rho = 8 \times 28.09/(6.022 \times 10^{23} \times
(5.431 \times 10^{-8})^3) = \qty{2.33}{g/cm^3}$，即实测值。
\end{solution}

\begin{solution}{exo:b1:crystals-ionic-covalent:12}
$x = 152/181 = 0.84 > 0.732$：该规则预言为氯化铯型。
若取氯化钠型结构，\ce{Rb-Cl} $= a/2 = \qty{329.1}{pm}$，接近
$152 + 181 = \qty{333}{pm}$：观测结构与之相符。该
规则把离子视为硬球，忽略了其极化率及
两种结构间能量的微小差别（氯化铷在高压下
取氯化铯型）；它只是一条指导性规律，
而非定律。
\end{solution}

\begin{solution}{pb:b1:crystals-ionic-covalent:1}
\textbf{1.} 位于 8 个顶点与 6 个面心处：$8 \times \frac18 + 6
\times \frac12 = 4$ 个离子。
\textbf{2.} 8 个四面体间隙全部被占据：8 个 \ce{F-} 对 4 个
\ce{Ca^{2+}}，比例为 1:2，即 \ce{CaF2}。
\textbf{3.} \ce{F-} 周围有 4 个 \ce{Ca^{2+}}，位于四面体的顶点。
\ce{Ca^{2+}} 周围有 8 个 \ce{F-}，位于立方体的顶点。
\textbf{4.} $a\sqrt3/4 = 546.3 \times 0.4330 = \qty{236.6}{pm}$.
\textbf{5.} $112 + 131 = \qty{243}{pm}$，误差在 3\,\% 以内。
\textbf{6.} \ce{F-} 离子构成简单立方排列，棱长为 $a/2 =
\qty{273.2}{pm}$，略大于 $2 \times 131 = \qty{262}{pm}$，两者
几乎相接触。
\textbf{7.} $x = 112/131 = 0.85$.
\textbf{8.} $x > 0.732$，配位数为 8。
\textbf{9.} 氯化铯型中阳离子与阴离子数目相等；若阴离子数是其两倍，
则只有半数的阴离子立方体能容纳阳离子。
\textbf{10.} 晶胞的 8 个 \ce{F-} 构成一个由 8 个小立方体组成的大立方体，
小立方体棱长为 $a/2$；\ce{Ca^{2+}} 交替占据其中 4 个的中心，
犹如三维棋盘上的黑格。
\textbf{11.} $C = \frac43\pi(4 \times 112^3 + 8 \times 131^3)/546.3^3 = 0.61$.
\textbf{12.} \ce{O^{2-}} 位于面心立方晶格（顶点与面心）上，
\ce{Li+} 位于 8 个四面体间隙中。
\textbf{13.} $a\sqrt3/4 = \qty{200.0}{pm}$; $59 + 142 = \qty{201}{pm}$.
\textbf{14.} $\rho = 4 \times 29.88/(6.022 \times 10^{23} \times (4.619
\times 10^{-8})^3) = \qty{2.01}{g/cm^3}$，即实测值。
\textbf{15.} 闪锌矿型：金刚石即是两类位置均由碳原子占据的
闪锌矿型结构。
\textbf{16.} $a\sqrt3/4 = \qty{154.5}{pm}$; $\rho = 8 \times
12.01/(6.022 \times 10^{23} \times (3.567 \times 10^{-8})^3) =
\qty{3.52}{g/cm^3}$.
\textbf{17.} 金刚石中只有半数四面体间隙被占据，占据者
与晶格原子等大，仅沿键相接触：
$C = 0.34$。萤石中所有间隙均被占据，较小的阳离子
填充在较大阴离子之间的空隙中：$C = 0.61$。
\textbf{18.} $40.08 + 2 \times 19.00 = \qty{78.08}{g/mol}$.
\textbf{19.} $4 \times 78.08/(6.022 \times 10^{23}) = \qty{5.186e-22}{g}$.
\textbf{20.} $(5.463 \times 10^{-8})^3 = \qty{1.630e-22}{cm^3}$.
\textbf{21.} 每一百个晶胞缺一个 \ce{F-}，所去质量为 $19.00/(100
\times 312.3)$ 的质量，即 0.06\,\%，在三位有效数字下不可见（且还须去掉一个
\ce{Ca^{2+}}，否则晶体不再呈电中性）。
\textbf{22.} $\rho = 5.186 \times 10^{-22}/1.630 \times 10^{-22} =
\textbf{\qty{3.18}{g/cm^3}}$，恰与实测密度相符：该模型
晶胞由一次 X 射线测得的 $a$ 构建，
可解释该矿物的质量。
\end{solution}
